\section{The projective case}\label{sec:projective}

We construct two complete directions from families of genus-one
curves. The curves on the K3 surface need not be smooth: the
complete flows take place on their smooth normalizations.

\begin{proposition}\label{prj:strips}
Let $S$ be a smooth projective complex K3 surface. There is a
nonempty Zariski open subset $U\subset S$ such that, for every
$x\in U$, there are $b>0$ and holomorphic maps
\[
\sigma_1:\C\times\Delta_b\longrightarrow S,\qquad
\sigma_2:\Delta_b\times\C\longrightarrow S
\]
which are locally biholomorphic at the origin and satisfy
\[
\sigma_1(0,0)=\sigma_2(0,0)=x,\qquad
\sigma_1(z,0)=\sigma_2(z,0)\quad (|z|<b).
\]
The complement $S\setminus U$ is algebraic of dimension at most
one. Consequently $S$ has property $L$ everywhere.
\end{proposition}

We use the genus-one case of the theorem of Chen and Gounelas
\cite[Theorem~A]{ChenGounelas}:
there are integral curves $C_n\subset S$ of geometric genus one
with $C_n^2\to\infty$, whose normalization maps belong to smooth
proper families of maps
\[
p_n:Y_n\longrightarrow T_n,\qquad q_n:Y_n\longrightarrow S.
\]
Here $T_n$ is an irreducible algebraic curve, the fibers of $p_n$
are connected genus-one curves, and the induced moduli map is
nonconstant. One fiber map is the normalization of $C_n$ followed
by its inclusion into $S$. A stable-map family includes, by definition,
a proper family of source curves; pulling back that family and restricting
to the smooth-source locus gives the proper maps $p_n$. No properness of
$T_n$ is asserted. We use the families over their full proper
fibers, even when evaluation is inverted only on a smaller open
subset.

\begin{lemma}\label{prj:prepare}
Let $S$ be a smooth projective integral surface, let $H$ be an
ample divisor, and let $p:Y\to T$ be a smooth proper family of
connected genus-one curves over a smooth irreducible algebraic
curve. Suppose that $q:Y\to S$ has nonconstant source-curve
moduli and that $q_{t_0}$ is the normalization map of an integral
curve $C\subset S$. After removing finitely many points of $T$,
every $q_t$ normalizes an integral image curve $C_t$ with
$H\cdot C_t=H\cdot C$. The evaluation map $q$ is dominant and
generically finite, and some nonempty Zariski open $V\subset S$
has finite etale surjective inverse image map $q^{-1}(V)\to V$.
\end{lemma}

\begin{proof}
The total space $Y$ is a smooth integral surface. Smoothness
follows from that of $p$ and $T$; connectedness follows from
connectedness of the base and fibers, and the irreducible
components of a smooth variety are disjoint open and closed
sets. Proper smooth local triviality shows that
$\deg(q_t^*H)$ is constant: it is the integral of $c_1(q^*H)$
on a fiber. At $t_0$ it equals $H\cdot C>0$. Thus all fiber
maps are nonconstant and finite onto their reduced integral
images.

Consider the graph map
\[
Q=(p,q):Y\longrightarrow T\times S.
\]
It is proper, by factoring through the closed graph in
$Y\times S$ followed by the proper map $p\times\mathrm{id}_S$,
and it has finite fibers. The proper quasi-finite criterion
therefore makes it finite. Its reduced image $Z$ is an integral
surface.
Choose a smooth point $z_0$ of $C$ and its unique preimage
$y_0\in Y_{t_0}$. The fiber $Q^{-1}(t_0,z_0)$ is a singleton.
Moreover $dQ_{y_0}$ is injective: a vector in its kernel is
vertical for $p$, and the normalization is an isomorphism near
$z_0$, so $dq$ is injective on that vertical tangent line.

The holomorphic constant-rank theorem embeds a small
neighborhood $W$ of $y_0$ in $T\times S$. Properness makes
$Q(Y\setminus W)$ closed in the analytic topology. Shrink a neighborhood
of $(t_0,z_0)$ to avoid that closed set and remain inside the
embedding chart. Thus $Q$ is an isomorphism over a nonempty
analytic open subset of $Z$. That subset meets the dense
Zariski open locus where the finite map is etale of its generic
degree. The degree is therefore one.

Since $Y$ is normal, $Q$ is the normalization of $Z$. Its
nonisomorphism locus $N\subset Z$ is closed of dimension at
most one. Remove the finitely many parameters of the vertical
curve components of $N$. Every horizontal component of $N$
has finite intersection with every remaining fiber, so the
irreducible support
$|Z_t|=q_t(Y_t)$ is not contained in $N$. The restriction $q_t$
is consequently finite and birational onto its reduced image
$C_t$, hence is its normalization. The projection formula gives
\[
H\cdot C_t=\deg(q_t^*H)=H\cdot C.
\]

If the closure of $q(Y)$ were a curve, each retained fiber map
would normalize that same integral curve. Their smooth source
curves would all be isomorphic, contrary to nonconstant moduli.
Thus $q$ dominates $S$. Source and target have dimension two,
so the induced function-field extension is finite and, in
characteristic zero, separable.

Here the desired finite etale locus must be justified even
though $q$ itself is not assumed proper. Choose affine opens
$V_0=\operatorname{Spec}B\subset Y$ and
$U_0=\operatorname{Spec}A\subset S$ with $q(V_0)\subset U_0$.
Every one of finitely many $A$-algebra generators of $B$ is
algebraic over $\operatorname{Frac}A$. Invert a common nonzero
element of $A$ to make their monic equations integral; the
resulting affine map is finite. Remove also the closure of
$q(Y\setminus V_0)$, of dimension at most one, and shrink to
the etale locus. This leaves a nonempty Zariski open $V$ on
which the full inverse image is the finite etale piece.
It is surjective because it is finite and dominant.
\end{proof}

Normalize the parameter curves supplied by Chen--Gounelas,
retaining a point above each specified normalization fiber,
and pull back the families. The bases are then smooth
irreducible curves; smoothness, properness, and nonconstant
moduli persist. Fix a very ample divisor $H$ on $S$. The
Hodge index inequality
\[
(H\cdot C_n)^2\ge H^2 C_n^2
\]
shows that the ample degrees are unbounded.
Choose two distinct degrees $d_1,d_2$ and apply
Lemma~\ref{prj:prepare}. We obtain two families
\[
p_i:Y_i\longrightarrow T_i,\qquad q_i:Y_i\longrightarrow S
\quad(i=1,2)
\]
whose fiber maps normalize integral curves of degree $d_i$.
Choose a common nonempty open $V\subset S$ on which
$Y_i^V=q_i^{-1}(V)\to V$ are both finite etale and surjective.
Only their inverse sheets are restricted to $V$: the maps
$q_i$ remain defined on all of $Y_i$.

\begin{lemma}\label{prj:transverse}
There is a nonempty Zariski open $U\subset V$ such that any
two lifts $y_i\in Y_i^V$ of a point $x\in U$ have distinct
transported vertical tangent lines
\[
dq_1\bigl(\ker(dp_1)_{y_1}\bigr)
\ne dq_2\bigl(\ker(dp_2)_{y_2}\bigr)
\quad\text{in }T_xS.
\]
The closure in $S$ of the exceptional locus has dimension
at most one.
\end{lemma}

\begin{proof}
Form the finite etale cover
\[
W=Y_1^V\times_VY_2^V\xrightarrow{\,Q\,}V,\qquad
\tau_i=p_i\circ\operatorname{pr}_i.
\]
Every irreducible component of the smooth surface $W$
dominates $V$, since a finite etale map is both open and closed.
The maps $\tau_i$ are submersions. Their algebraic rank-drop
locus
\[
I=\{d\tau_1\wedge d\tau_2=0\}
\]
records equality of the transported vertical directions.

No component of $W$ can be contained in $I$. If one were,
identify a small analytic neighborhood in it with an open
subset of $S$ using $Q$. In a coordinate $s=\tau_1$,
dependence of the differentials makes $\tau_2$ depend only on
$s$. A local $\tau_1$-fiber would then lie in a local
$\tau_2$-fiber. Its nonconstant image germ in $S$ lies in an
integral curve from each family. Two distinct integral curves
on a projective surface have finite intersection, so the
curves must coincide. This contradicts their distinct
$H$-degrees.

Consequently $\dim I\le1$. Since $Q$ is finite, $Q(I)$ is
closed in $V$ and has dimension at most one; its algebraic
closure in $S$ has the same dimension bound. Taking
\[
U=S\setminus\bigl((S\setminus V)\cup\overline{Q(I)}^{\,S}\bigr)
\]
proves all assertions.
\end{proof}

\begin{lemma}[Complete vertical flows]\label{prj:flow}
Let $p:Y\to\Delta$ be a proper holomorphic submersion with
connected genus-one fibers. After shrinking the base disc,
there is a nowhere-zero holomorphic vertical vector field
whose flow is jointly holomorphic for all complex times.
Its restriction to the central fiber may be prescribed as
any nonzero holomorphic vector field there.
\end{lemma}

\begin{proof}
For every fiber $C$, the space $H^0(C,T_C)$ is one-dimensional
and its nonzero elements have no zeros. The relative tangent
bundle is locally free and flat over the base. Grauert's
cohomology and base-change theorem \cite[\S~7, Satz~5]{Grauert} makes
$p_*T_{Y/\Delta}$ a holomorphic line bundle with these spaces
as fibers. The evaluation morphism
\[
p^*(p_*T_{Y/\Delta})\longrightarrow T_{Y/\Delta}
\]
is an isomorphism on every fiber, hence everywhere.
A local frame gives the vertical field, and scaling it
prescribes its central restriction.

Take a smaller closed base disc. Its full inverse image is
compact. Local holomorphic flow existence and a finite cover
of this compact set give a common time radius $\epsilon>0$.
The flow preserves fibers, so successive small-time flows
remain in the same compact inverse image. Composing the
time-$z/M$ flow $M$ times defines a holomorphic flow for
$|z|<M\epsilon$. Uniqueness and the local group law make
these extensions agree as $M$ increases. They define the
jointly holomorphic flow for every $z\in\C$ over the smaller
open base disc.
\end{proof}

\begin{proof}[Proof of Proposition~\ref{prj:strips}]
Fix $x\in U$ and lifts $y_i$ as in Lemma~\ref{prj:transverse}.
Use Lemma~\ref{prj:flow} on each full proper family over a
small base disc about $p_i(y_i)$. Let $\xi_i$ be the resulting
vertical fields and $\operatorname{Fl}_i^z$ their complete
flows. Then
$v_i=dq_i(\xi_i(y_i))$ are independent tangent vectors at $x$.

Choose a local section $e_1(w)$ of $p_1$ through $y_1$ and set
\[
\sigma_1(z,w)=q_1\bigl(\operatorname{Fl}_1^z(e_1(w))\bigr).
\]
This map is defined for all $z\in\C$ and all sufficiently
small $w$. It is locally biholomorphic at the origin:
the flow derivative is vertical, the section derivative
projects nontrivially to the base, and $q_1$ is etale at $y_1$.
Put $\varphi(z)=\sigma_1(z,0)$.

For small $z$ lift $\varphi(z)$ through the local inverse
of $q_2$ at $y_2$, obtaining a holomorphic map $\ell_2(z)$
with $q_2(\ell_2(z))=\varphi(z)$. Shrink its domain so
$p_2(\ell_2(z))$ lies in the disc supporting the second
complete flow, and define
\[
\sigma_2(z,w)=q_2\bigl(\operatorname{Fl}_2^w(\ell_2(z))\bigr).
\]
It is defined for small $z$ and every $w\in\C$.
Its derivatives at the origin are $v_1,v_2$, so it too
is locally biholomorphic. At time zero it returns its
starting point, giving the exact identity
$\sigma_2(z,0)=\varphi(z)=\sigma_1(z,0)$.
Choose a common radius for the bounded variables.
The inverse sheet was used only for the initial lift;
subsequent flow values lie in entire proper fibers,
where $q_2$ is still defined.

Proposition~\ref{loc:criterion} gives $L$ throughout $U$.
Its algebraic complement has dimension at most one, so
Corollary~\ref{loc:disc} gives $L$ at the remaining points.
\end{proof}

